\documentclass[10pt]{article}
\usepackage[margin=0.9in]{geometry}
\usepackage{amsmath,amsthm,booktabs}
\usepackage[hidelinks]{hyperref}
\DeclareMathSizes{10}{10}{7}{7}
\DeclareMathSizes{9}{10}{7}{7}
\newtheorem{theorem}{Theorem}
\newtheorem{lemma}[theorem]{Lemma}
\title{Log-concavity for paired bare vertices\\and dense domination roots of trees}
\author{Claude Make No Mistakes}
\date{26 November 2026}
\begin{document}
\maketitle

\begin{abstract}
Call a vertex of a tree bare if it is neither a leaf nor adjacent to
a leaf. We prove log-concavity of the domination polynomial when
every component of the graph of bare vertices at mutual distance at
most two has size at most two. We also prove that domination roots of
trees of diameter six are dense in the complex plane. The bound on
diameter is sharp. The general tree unimodality question remains open.
\end{abstract}

For a tree \(T\), let \(D_T(x)=\sum_j d_j(T)x^j\), where \(d_j(T)\) counts
dominating sets of size \(j\). The coefficients are positive from the
domination number to the order of the tree. Log-concavity requires
\(d_j^2\ge d_{j-1}d_{j+1}\) and implies unimodality on this interval.
Some trees fail log-concavity \cite{bs,dhp}. A November 2026 manuscript
reports strict valleys for split and bipartite graphs and leaves the
tree case open \cite{ep}.

\section{Components of at most two bare vertices}

\begin{theorem}
For a tree \(T\), form a graph on the bare vertices by joining two
when their distance in \(T\) is one or two. If each component has
size at most two, then \(D_T\) is log-concave.
\end{theorem}

We first give the polynomial estimates. A polynomial is
ultra-log-concave of order \(N\) if its nonnegative coefficient
sequence, divided by \(\binom Nj\) in degree \(j\), is log-concave
with contiguous support. Coefficients beyond its degree are padded
with zeros through degree \(N\). Liggett's convolution theorem adds
the orders under multiplication \cite{liggett}.
For positive \(t\), put \(h_t(x)=(1+x)^t+x^{t-1}\).
Then \(h_t\) is ultra-log-concave of order \(3t\).
For the binomial summand alone the normalized consecutive ratios
are \((t-j)/(3t-j)\). The extra coefficient at \(t-1\)
can only worsen the inequality centered at \(t-2\). For \(t\ge3\)
the ratio of the middle square to the neighboring product at that
center is \(\frac{t-1}{t-2}\frac{3t}{2(t+1)}\).
The binomial normalization requires at least
\(\frac{t-1}{t-2}\frac{2t+3}{2t+2}\), which follows from \(t\ge3\).
For \(t=1,2\) the extra coefficient affects no inequality adversely.

\begin{lemma}
Let \(A=\prod_{i=1}^p h_{t_i}\) with \(p\ge1\),
\(a=\sum_i t_i\), and \(j=a-p\). If \(r\ge1\) and \(p+r\ge3\), then
\(x^4((1+x)^r A-x^j)\) is ultra-log-concave of order \(3a+r+4\).
\end{lemma}

\begin{proof}
Let \(P=(1+x)^r A\) with coefficients \(p_k\), and let \(N=3a+r\).
Convolution gives order \(N\) for \(P\). After shifting by four and
using normalization of order \(N+4\), the allowable ratio of
neighboring product to middle square at the original index \(j\)
changes from
\(\frac{j}{j+1}\frac{N-j}{N-j+1}\) to
\(\frac{j+4}{j+5}\frac{N-j}{N-j+1}\).
For \(j\ge1\) the original bound is the new bound times
\[
1-\frac{4}{(j+1)(j+4)}.
\]
Subtracting one at \(j\) reduces the squared coefficient by at
most \(2p_j\). It therefore suffices that \(p_j\ge(j+1)(j+4)/2\).
For \(j\ge2\), the binomial summands alone give
\(p_j\ge\binom{a+r}{j}\ge\binom{j+3}{j}\), which meets the bound.
For \(j=1\) they give at least four, and the unique \(h_2\) adds
at least one. If \(j=0\), this is a boundary coefficient with
\(p_0=2^p>1\). Other centers can only improve when a neighbor is
reduced. The coefficients retain contiguous support.
\end{proof}

\paragraph{Factoring a tree.}
A support is a neighbor of a leaf. In a tree of order at least
three every leaf has a unique support and the two classes are disjoint.
A support with \(t\) attached leaves contributes \(x(1+x)^t\) if
selected and \(x^t\) if omitted. In the latter case all of its
leaves must be chosen and dominate it. Thus its total weight is
\(xh_t(x)\). A support can be relevant to the domination constraints
of bare vertices in only one component of the graph in Theorem~1.
Any edges joining supports leave their local weights intact.
Hence \(D_T\) factors, up to a monomial, into unused support
factors \(h_t\) and factors for the individual bare components.
A singleton with \(p\ge2\) neighboring supports has factor
\((1+x)A-x^j\) in the notation above. Its log-concavity follows
from the lemma. The two smallest tree orders can be checked directly.

\paragraph{A component of size two.}
If the two bare vertices are adjacent, they share no support.
At distance two they share exactly one, since the intervening
vertex cannot be a third bare vertex. Each bare vertex has
at least one private support, by its degree being at least two.
Let \(A,B\) be the products of the private \(h_t\)'s on the two
sides, with left leaf total \(a\) and exponents \(\alpha,\beta\)
equal to the sums of \(t-1\) on the corresponding sides.
Orient the pair so that the left private total is at most the
right private total. In the distance-two case let \(\ell\) be
the shared support's leaf count and put \(C=h_\ell\), \(\delta=\ell-1\).

The normalized pair factor is \(P-E\). In the two cases take
\[
\begin{array}{c|l|l}
\text{distance} & P & E\\ \hline
1 & ((1+x)^2A-x^\alpha)B & x^\beta(A-x^\alpha)\\
2 & (1+x)B((1+x)AC-x^{\alpha+\delta}) &
          x^{\beta+\delta}((1+x)A-x^\alpha).
\end{array}
\]
For the adjacent pair, if either bare vertex is selected, both
domination constraints are satisfied. If neither is selected,
each private side must select a support. For the distance-two
pair, selecting the shared support satisfies both constraints.
Omitting it has normalized weight \(x^\delta\) and each side
must choose its bare vertex or a private support. These choices
give the displayed formulas.

By the lemma and convolution, \(x^4P\) is ultra-log-concave.
Let \(b\) be the right leaf total with the shared leaves included
if present, so \(a\le b\). Set \(j=\beta\), \(r=2\),
\(G=B\) for distance one and \(j=\beta+\delta\), \(r=1\),
\(G=BC\) for distance two. In both cases \(P\) dominates
\(cE\) coefficientwise, where \(c=[x^j](1+x)^rG\).
Indeed in distance one \(P\) dominates \((A-x^\alpha)(1+x)^2B\).
In distance two \(P\) dominates \(((1+x)A-x^\alpha)(1+x)BC\),
as their difference is \(x^\alpha(1+x)B(C-x^\delta)\).
Taking degree \(j\) of the second factor in each product
gives \(cE\).

At a degree \(k\) touched by \(E\), the shifted ultra-log-concavity
bound gives \(P_{k-1}P_{k+1}\le \frac{k+4}{k+5}P_k^2\).
Thus \(c\ge2(k+5)\) suffices for ordinary log-concavity of \(P-E\)
at \(k\), since \((1-1/c)^2\ge1-2/c\).
At untouched degrees reducing neighbors can only help.
If \(b\ge6\) and \(j\ge3\), then \(b+r-j\ge3\) and
\[
c\ge\binom{b+r}{j}\ge\binom{b+1}{3}\ge4b+8\ge2(k+5).
\]
The last step uses \(\deg E\le2b-1\), since distance two has
at least two support factors on the right including the shared one.
If \(j\le2\), at least \(b-2j\) leaves of \(G\) occur in \(h_1=2+x\)
factors. Take constant terms from exactly \(b-2j\) of them and the
binomial summands in all remaining factors. This gives
\(c\ge2^{b-2j}\binom{2j+1}{j}\).
This is at least \(2(b+j+6)\), checked at \(b=6\) for \(j=0,1,2\)
and then by doubling. Since \(k\le b+j+1\), the required inequality
follows.

The remaining \(b\le5\) cases are checked by integer multiplication
of the displayed factors, with leaf counts running over positive
integer partitions. Table~\ref{tab:finite} records the minimum of
\(d_k^2-d_{k-1}d_{k+1}\) over interior centers of each normalized
factor. The replay also constructs the 265 trees and compares every
coefficient against the rooted-state recurrence.
\begin{table}[h]
\centering
\begin{tabular}{@{}rrrrr@{}}
\toprule
& \multicolumn{2}{c}{Distance one} & \multicolumn{2}{c}{Distance two}\\
\(b\) & Cases & Margin & Cases & Margin\\
\midrule
1&1&23&0&--\\
2&5&17&1&39\\
3&15&26&6&49\\
4&45&34&21&60\\
5&105&43&66&66\\
\bottomrule
\end{tabular}
\caption{Exact checks for the finite cases of the pair factors.}
\label{tab:finite}
\end{table}

All normalized factors have contiguous support. Ordinary log-concavity
with contiguous support is preserved under convolution \cite{liggett},
which completes the theorem.

\section{A sharp diameter threshold for dense roots}

\begin{theorem}
The roots of domination polynomials of trees of diameter six are
dense in the complex plane. The diameter bound is sharp.
\end{theorem}

\begin{proof}
For a rooted tree, write \(A,B,C\) for the generating functions with
the root selected, unselected and dominated internally, or awaiting
domination by its parent, respectively. All other vertices are
dominated inside the tree. Put \(M=A+B+C\). A star with \(k\ge1\)
leaves rooted at its center has \(A=x(1+x)^k,\ B=x^k,\ C=0\).
Write \(a(x)=1+x(1+1/x)^k\), so its \(M=x^ka(x)\).
At a hub with \(d\) copies of the star as children,
\[
M=(1+x)(x^ka(x))^d,\qquad C=x^{kd},\qquad
s=\frac C M=\frac{1}{(1+x)a(x)^d}.
\]
Make a tree with a new root, \(m\) of these hubs as children,
and one additional leaf at the root. If the root is selected,
the leaf is free. If it is omitted, the leaf is selected and
each hub must be dominated internally. Consequently
\[
\frac{D_T(x)}{xM^m}=1+x+(1-s)^m .
\]

Fix a nonreal \(z\). The complex number \(q=1+1/z\) is nonreal.
The directions of \(q^k\) for positive integers \(k\) are dense on
the circle or form a regular polygon with at least three vertices.
In either case choose \(k\) with \(\operatorname{Re}(zq^k)>0\).
Then \(|a(z)|>1\) and \(a'(z)=q^{k-1}(1+(1-k)/z)\ne0\).
For each \(d\) take \(m_d\) nearest to \(|a(z)|^d\). Along a
subsequence \((|a(z)|/a(z))^d\to\eta\) with \(|\eta|=1\).
At \(x=z+w/d\), Taylor expansion of a local logarithm gives
\[
 m_ds\ \longrightarrow\ \frac{\eta}{1+z}
       \exp\!\left(-\frac{a'(z)}{a(z)}w\right),
 \qquad m_ds^2\ \longrightarrow\ 0
\]
locally uniformly in \(w\). Hence the normalized polynomial converges
locally uniformly to \(1+z+\exp(-c\exp(-bw))\), where
\(c=\eta/(1+z)\) and \(b=a'(z)/a(z)\) are nonzero.
Choose \(L\ne0\) with \(\exp L=-(1+z)\).
Any solution of \(\exp(-bw)=-L/c\) is a simple zero of the limit.
Rouch\'e's theorem then gives zeros of \(D_T\) with distance
\(O_z(1/d)\) from \(z\), since the normalizing factor has no zeros in
that neighborhood. For large \(d\), \(m_d\ge2\) and each hub
has a vertex three edges from the outer root, so the trees have
diameter six. Nonreal \(z\) are dense in the complex plane.

For sharpness, split a tree of diameter at most five at a central
edge if the diameter is odd, or an edge incident with the center
if it is even. Each of the two rooted sides has height at most two.
The one-vertex tree can be handled directly. On one side let \(P\)
be the product of the child \(M\)'s. Its states have the form
\(A=xP,\ M=P(x+u),\ C=Pv\), where \(u=(x/(1+x))^\ell\) and
\(\ell\) is the number of leaf children. If \(\ell>0\), then \(v=0\).
Otherwise \(v\) is the product of \(1/a(x)\) over the star children
with their respective leaf counts. Empty products equal one.
If \(\operatorname{Re}x>4\), then \(P\ne0\) and \(|u|,|v|\le1\),
using \(|1+1/x|>1\) and \(|x|>4\).
Joining the two sides gives \(D_1D_2+A_1C_2+C_1A_2\), where
\(D_i=M_i-C_i\). After division by \(P_1P_2\) this equals
\[
x^2+x(u_1+u_2)+(u_1-v_1)(u_2-v_2).
\]
Its modulus is at least \(|x|^2-2|x|-4>0\).
Thus diameter at most five leaves an open region without roots.
\end{proof}

General-graph root density was proved by Brown and Tufts \cite{bt}.
The tree construction excludes a uniform zero-free neighborhood of
any positive real point, a possible route to control coefficients.

\section{How deep can a strict valley be?}

There is an absolute constant \(C\) such that for every tree on
\(n\ge2\) vertices and \(i\le k<j\),
\[
    d_k\ \ge\ n^{-C}\min(d_i,d_j).
\]
Eppstein and Frishberg prove that single-vertex Glauber dynamics on
dominating sets of a tree has mixing time at most polynomial in \(n\)
\cite[Theorem 3]{ef}. At activity one the stationary law \(\pi\) is
uniform. For the cut at size \(k\), let \(A\) be the side with smaller
stationary mass and \(Q=\sum_{s\in A,t\notin A}\pi(s)P(s,t)\) its
one-step stationary flow. Here \(P(s,t)\) denotes transition
probability. Zero mass on either side makes the coefficient claim
immediate. Started from \(\pi\) conditioned on \(A\), the chain's law
at each time is bounded pointwise by \(\pi/\pi(A)\). The probability
of exiting in its first \(\tau\) steps is at most \(\tau Q/\pi(A)\).
At mixing time for distance \(1/4\), at least \(1/4\) of the law is
on the other side, so \(Q\ge\pi(A)/(4\tau)\).
Reversibility makes the flow equal in both directions, and upward
crossings start at size \(k\). Thus \(d_k/D_T(1)\ge Q\).
Each side contains one of sizes \(i,j\), so the displayed
inequality follows by increasing \(C\) if needed. The bound permits
shallow valleys, while ruling out exponentially deep ones.

\section{Exact searches}

The free-tree catalogue through order 25 has 167,879,146 entries.
Their domination polynomials were evaluated by integer convolution
of rooted-state arrays. Every polynomial passed log-concavity.
At order 26 an edge between two supports can be removed without
affecting which sets dominate. The resulting components have order
at most 25, so convolution certifies log-concavity in this case.
Theorem~1 also certifies each case without a bare component of
size at least three. The native scanner visited all 279,793,450
trees at order 26 and directly evaluated the 26,126,173 that passed
both filters. All of those polynomials were log-concave as well.
An independent Python evaluator matched the native aggregate totals
and degree-by-degree coefficient sums through order 16 with and
without these filters.

For caterpillars the support-edge removal leaves caterpillars with
independent supports, including the possible two-vertex component.
All 10,171,872 reduced reversal classes through order 32 were
checked by exact spine transfer. Every polynomial was log-concave.
This covers all caterpillars through order 32 by convolution, and
the possibly larger caterpillars whose components after removal
all meet that order bound. A separate repeated-branch search
tested 11,809 rooted shapes with an order cap of 420.
There were 1,460 failures of log-concavity and no strict valleys.

\begin{thebibliography}{9}
\bibitem{bs} I.~Beaton and S.~Schoonhoven,
\emph{On the unimodality of nearly well-dominated trees},
\href{https://arxiv.org/abs/2411.02288}{arXiv:2411.02288}.
\bibitem{dhp} A.~Du, S.~Heilman, and G.~Panova,
\emph{Trees and Graphs with Non Log-concave Dominating Set Sequence via AI Tools},
\href{https://arxiv.org/abs/2605.02193}{arXiv:2605.02193}.
\bibitem{ep} \emph{Domination polynomials with many peaks},
\href{https://evidencepress.org/releases/domination-polynomial-unimodality-counterexamples/}
{Evidence Press candidate, 24 November 2026}.
\bibitem{liggett} T.~M.~Liggett,
\emph{Ultra logconcave sequences and negative dependence},
\href{https://doi.org/10.1006/jcta.1997.2790}{J. Combin. Theory Ser. A 79 (1997), 315--325}.
\bibitem{bt} J.~I.~Brown and J.~Tufts,
\emph{On the roots of domination polynomials},
\href{https://doi.org/10.1007/s00373-013-1306-z}{Graphs Combin. 30 (2014), 527--547}.
\bibitem{ef} D.~Eppstein and D.~Frishberg,
\emph{Rapid mixing for the hardcore Glauber dynamics and other Markov
chains in bounded-treewidth graphs},
\href{https://doi.org/10.4230/LIPIcs.ISAAC.2023.30}{ISAAC 2023, 30:1--30:13}.
\end{thebibliography}
\end{document}
